% Unit 6 free-response set -- 2 long (10) + 3 short (4) = 32 pts.
\documentclass{apchem-exam}

\apcunit{6}
\apcunittitle{Thermochemistry}
\apcdoctitle{Free-Response Set}
\apcsubtitle{CED 6.1--6.9 \textbullet\ 5 questions \textbullet\ 32 points \textbullet\ 50 minutes}

\begin{document}
\apctitleblock

\begin{apcnote}[Scope]
Covers \ced{6.1} through \ced{6.9}, worth
\SI{7}{\percent}--\SI{9}{\percent} of the exam. Note that \ced{6.5}
(energy of phase changes) is sourced from Zumdahl \S10.8 and \ced{6.7}
(bond enthalpies) from \S8.8 --- both sit outside the thermochemistry
chapter, and both are examined here.

Topics \ced{6.6} and \ced{6.9} carry an exclusion covering the
\textbf{enthalpy-versus-internal-energy distinction} and formal treatment
of \textbf{state functions}. Hess's law itself is fully examinable; the
thermodynamic formalism behind it is not.

Take $c(\ce{H2O}) = \SI{4.18}{\joule\per\gram\per\celsius}$,
$\Delta H_{\text{fus}} = \SI{334}{\joule\per\gram}$ and
$\Delta H_{\text{vap}} = \SI{2260}{\joule\per\gram}$.
\end{apcnote}

\examsection{II}{Free Response}{5 questions \textbullet\ 32 points \textbullet\ 50 minutes}{\calculatorok}

% =====================================================================
\frq{long}{10}{\ced{6.3} \ced{6.4} \ced{6.6} \textbullet\ calorimetry}

A \SI{75.0}{\gram} block of an unknown metal is heated to
\SI{98.0}{\celsius} and dropped into \SI{150.0}{\gram} of water at
\SI{22.0}{\celsius} in an insulated calorimeter. The final temperature of
both is \SI{24.6}{\celsius}.

\begin{parts}
  \item State the direction of heat flow and the sign of $q$ for the water
        and for the metal. \answerlines[3]{Heat flows from the
        \textbf{metal} (hotter) to the \textbf{water} (cooler), until both
        reach the same temperature. $q_{\text{water}}$ is
        \textbf{positive} --- it gains energy; $q_{\text{metal}}$ is
        \textbf{negative} --- it loses the same amount.}
  \item Calculate the heat gained by the water. \answerlines[2]{$q = mc
        \Delta T = (150.0)(4.18)(24.6 - 22.0) = \mathbf{1.63\times10^{3}}$~J}
  \item Calculate the specific heat capacity of the metal.
        \workspace[3.2cm]
        \keyonly{\begin{apcanswer}The metal lost what the water gained,
        cooling from \SI{98.0}{\celsius} to \SI{24.6}{\celsius} --- a fall
        of \SI{73.4}{\celsius}, \emph{not} \SI{2.6}{\celsius}:

        $c = \dfrac{1630}{(75.0)(73.4)} =
        \mathbf{0.296}$~J/g\,\si{\celsius}\end{apcanswer}}
  \item The water's temperature rose only \SI{2.6}{\celsius} while the
        metal fell \SI{73.4}{\celsius}, for the same quantity of heat.
        Explain. \answerlines[3]{Water has a much larger specific heat
        capacity --- about fourteen times the metal's --- and there is
        twice as much of it by mass. The same energy therefore produces a
        far smaller temperature change in the water. Its extensive heat
        capacity, $mc$, is much the larger of the two.}
  \item State one assumption made in this calculation.
        \answerlines[2]{That the calorimeter is perfectly insulated, so
        all the heat lost by the metal is gained by the water and none
        escapes to the surroundings or is absorbed by the container.}
\end{parts}

\begin{rubric}
  \rpoint{2}{(a) 1 pt direction of flow; 1 pt both signs}
  \rpoint{2}{(b) 1 pt $mc\Delta T$ with $\Delta T = \SI{2.6}{}$; 1 pt
    \SI{1.63e3}{\joule}}
  \rpoint{3}{(c) 1 pt equating heat lost to heat gained; 1 pt using the
    metal's own $\Delta T$ of \SI{73.4}{\celsius}; 1 pt
    \num{0.296}~J/g\,\si{\celsius}. Using \SI{2.6}{\celsius} for the metal
    caps this part at 1}
  \rpoint{2}{(d) 1 pt water's larger specific heat capacity; 1 pt the
    larger mass / larger $mc$}
  \rpoint{1}{(e) any valid isolation assumption}
\end{rubric}

% =====================================================================
\frq{long}{10}{\ced{6.7} \ced{6.8} \ced{6.9} \textbullet\ Hess's law, formation and bond enthalpies}

Consider these data at \SI{298}{\kelvin}:
\begin{center}
\begin{tabular}{lr}
\hline
\textbf{Reaction} & $\Delta H^\circ$ (kJ) \\
\hline
\ce{N2(g) + O2(g) -> 2NO(g)}    & $+180.5$ \\
\ce{2NO(g) + O2(g) -> 2NO2(g)}  & $-114.1$ \\
\hline
\end{tabular}
\end{center}

\begin{parts}
  \item Use Hess's law to find $\Delta H^\circ$ for
        \ce{N2(g) + 2O2(g) -> 2NO2(g)}. \workspace[2.6cm]
        \keyonly{\begin{apcanswer}The two equations add directly --- the
        \ce{2NO} produced by the first is consumed by the second and
        cancels, and the oxygens sum to \ce{2O2}:

        $\Delta H^\circ = (+180.5) + (-114.1) =
        \mathbf{+66.4}$~kJ\end{apcanswer}}
  \item Hence determine $\Delta H_f^\circ$ for \ce{NO2(g)}, and explain
        the relationship. \answerlines[3]{The reaction in (a) forms
        \textbf{two} moles of \ce{NO2} from its elements in their standard
        states, so it equals $2\,\Delta H_f^\circ(\ce{NO2})$.

        $\Delta H_f^\circ = 66.4/2 = \mathbf{+33.2}$~kJ/mol}
  \item State $\Delta H^\circ$ for \ce{2NO2(g) -> 2NO(g) + O2(g)} and
        justify. \answerlines[2]{$+114.1$~kJ. Reversing a reaction
        reverses the direction of energy flow, so the magnitude is
        unchanged and the sign flips.}
  \item Using average bond enthalpies \ce{N#N} 941, \ce{O=O} 495 and
        \ce{N=O} 607 \si{\kilo\joule\per\mole}, estimate $\Delta H$ for
        \ce{N2 + O2 -> 2NO} and compare with the tabulated $+180.5$~kJ.
        \workspace[3cm]
        \keyonly{\begin{apcanswer}Broken: $941 + 495 = \num{1436}$~kJ.
        Formed: $2(607) = \num{1214}$~kJ.

        $\Delta H = 1436 - 1214 = \mathbf{+222}$~kJ

        That is about \SI{40}{\kilo\joule} above the tabulated
        $+180.5$~kJ. Bond enthalpies are \emph{averages} taken across many
        compounds, so a value computed this way is only an estimate; the
        sign and rough magnitude are right, the precise figure is
        not.\end{apcanswer}}
\end{parts}

\begin{rubric}
  \rpoint{2}{(a) 1 pt the equations add with \ce{NO} cancelling; 1 pt
    $+66.4$~kJ}
  \rpoint{3}{(b) 1 pt recognizing the reaction forms 2 mol from elements
    in standard states; 1 pt dividing by 2; 1 pt $+33.2$~kJ/mol}
  \rpoint{2}{(c) 1 pt $+114.1$~kJ; 1 pt the sign-reversal reason}
  \rpoint{3}{(d) 1 pt bonds broken; 1 pt bonds formed and the subtraction
    in the correct sense; 1 pt noting these are averages, so the result is
    an estimate. Reversing the subtraction gives $-222$ and caps this at 1}
\end{rubric}

% =====================================================================
\frq{short}{4}{\ced{6.1} \ced{6.2} \textbullet\ endothermic processes and energy diagrams}

Barium hydroxide octahydrate and ammonium chloride are mixed in a beaker.
The beaker becomes so cold that it freezes to a wet block of wood.

\begin{parts}
  \item State the sign of $\Delta H$ and explain the direction of heat
        flow. \answerlines[3]{$\Delta H$ is \textbf{positive} --- the
        process is endothermic. Energy flows \emph{from} the surroundings
        (the beaker, the wood, the air) \emph{into} the reacting mixture,
        so the surroundings lose energy and their temperature falls.}
  \item Describe the energy diagram for this reaction.
        \answerlines[2]{The products lie \textbf{above} the reactants,
        since the system absorbed energy. The vertical gap between them is
        $\Delta H$, and it is positive.}
\end{parts}

\begin{rubric}
  \rpoint{2}{(a) 1 pt positive/endothermic; 1 pt heat flows from
    surroundings into the system. ``It releases cold'' earns 0}
  \rpoint{2}{(b) 1 pt products above reactants; 1 pt the gap identified as
    $\Delta H$}
\end{rubric}

% =====================================================================
\frq{short}{4}{\ced{6.5} \textbullet\ energy of phase changes}

\begin{parts}
  \item Calculate the total energy needed to convert \SI{25.0}{\gram} of
        ice at \SI{0}{\celsius} into steam at \SI{100}{\celsius}.
        \workspace[3.2cm]
        \keyonly{\begin{apcanswer}Melt: $(25.0)(334) = \SI{8350}{\joule} =
        \SI{8.35}{\kilo\joule}$

        Heat \SI{0}{} to \SI{100}{\celsius}: $(25.0)(4.18)(100) =
        \SI{10450}{\joule} = \SI{10.45}{\kilo\joule}$

        Vaporize: $(25.0)(2260) = \SI{56500}{\joule} =
        \SI{56.5}{\kilo\joule}$

        Total $= \mathbf{75.3}$~kJ\end{apcanswer}}
  \item State which step requires the most energy and why.
        \answerlines[2]{\textbf{Vaporization.} Melting only loosens the
        lattice --- the molecules stay in contact --- whereas vaporizing
        must separate them completely, overcoming essentially all the
        remaining intermolecular attractions.}
\end{parts}

\begin{rubric}
  \rpoint{3}{(a) 1 pt each of the three steps computed correctly, with the
    phase-change terms using mass only and no $\Delta T$}
  \rpoint{1}{(b) vaporization, justified by having to separate the
    molecules completely}
\end{rubric}

% =====================================================================
\frq{short}{4}{\ced{6.3} \textbullet\ heat transfer and thermal equilibrium}

\begin{parts}
  \item Two blocks at different temperatures are placed in contact and
        isolated. Describe what happens at the particle level until they
        reach thermal equilibrium. \answerlines[3]{Particles of the hotter
        block have the greater average kinetic energy. At the boundary
        they collide with particles of the cooler block and transfer
        energy to them. The hotter block's average kinetic energy falls
        and the cooler block's rises, until both have the same average
        kinetic energy --- the same temperature. Collisions continue, but
        with no further \emph{net} transfer.}
  \item State whether the two blocks must contain the same amount of
        thermal energy at equilibrium, and why.
        \answerlines[2]{No. They share the same \emph{temperature}, but
        total thermal energy also depends on mass and specific heat
        capacity, which may differ between the blocks.}
\end{parts}

\begin{rubric}
  \rpoint{2}{(a) 1 pt energy transferred by particle collisions at the
    boundary; 1 pt equal average kinetic energy defines equilibrium}
  \rpoint{2}{(b) 1 pt ``no''; 1 pt temperature is not the same thing as
    total thermal energy}
\end{rubric}

\examtotal

\end{document}
